\label{sec:algorithm}

\subsection{Calibration Approach}

The default weights in Table~\ref{tab:signals} were not assigned arbitrarily.
They reflect a calibration argument that balances three considerations:
(1) legal recognition of the community type, (2) data quality and coverage,
and (3) empirical effect size on known test cases.
This section justifies each weight tier.

\subsection{Zone Co-membership: 0.30 (Highest)}

M.6 zone co-membership receives the highest weight because administrative
zone co-membership is the most consistently and explicitly recognized
community type in redistricting law.

\paragraph{Legal recognition.}
School district co-membership is cited as a redistricting criterion in
California Proposition 11 \citeyearpar{caprop11}, Colorado Article V
Section 44 \citeyearpar{coart5s44}, Arizona Proposition 106
\citeyearpar{azircprop106}, and Washington RCW 44.05.090 \citeyearpar{warcw44}.
In these states, splitting a school district across congressional districts
is a recognized harm that must be justified.
No other M-track signal has this direct statutory expression.

\paragraph{Mechanism.}
School district co-membership is a binary signal: two tracts are either
in the same school district or they are not.
This binary character makes the legal argument precise: if the plan splits
school district $X$, the expert can identify the specific tracts that were
separated and the specific school district boundary that was crossed.
The argument is not ``these tracts have similar economic character''; it is
``these tracts are governed by the same school board and their residents vote
in the same school board elections.''

\paragraph{Coverage.}
TIGER/Line school district boundaries are available for all 50 states and
cover virtually all census tracts.
This universal availability justifies placing M.6 in the minimum signal set.

\subsection{Economic Character and Land Use: 0.20 Each (Tied)}

M.1 (economic character) and M.2 (land use) share the second tier at 0.20 each.

\paragraph{M.1 justification.}
Economic character receives 0.20 because:
\begin{itemize}
  \item It has the clearest empirical effect among all M-track signals
    (14.3\,pp NC-14 proportionality gap improvement in the M.1/B.27 single-run
    result).
  \item Economic communities of interest are explicitly recognized in five
    or more state redistricting criteria.
  \item LODES WAC data has the highest data quality among all M-track sources:
    near-complete UI coverage, no sampling error, annual release.
\end{itemize}

\paragraph{M.2 justification.}
Land use receives 0.20 because:
\begin{itemize}
  \item Physical land use is the oldest recognized geographic community signal
    in redistricting, predating the modern communities-of-interest doctrine.
  \item The hard-cut rule for water boundaries provides a legally enforceable
    constraint not available in any other M-track signal.
  \item NLCD data is publicly available, CONUS-complete, and produced by
    the USGS with documented classification accuracy.
\end{itemize}

The parity between M.1 and M.2 reflects the judgment that economic and
physical community signals are equally important in the absence of a state-specific
argument for weighting one over the other.
Practitioners in states with strong land use transition evidence (coastal states,
states with clear urban-rural boundaries) may want to increase M.2's weight;
practitioners in states with concentrated employment clusters may want to
increase M.1's weight.

\subsection{Housing Character: 0.15}

M.3 receives a somewhat lower weight than economic character and land use
for two reasons.
First, housing character is slightly more correlated with race and income
than land use or employment data, creating a marginally higher legal exposure
in states with VRA-sensitive redistricting.
Second, the housing signal is partially redundant with M.2 (land use residential
classification and ACS housing density correlate), so some information is
double-counted.
The 0.15 weight captures the genuine additional information in M.3 (housing
tenure, age of housing stock, rental fraction) without over-weighting signals
that are correlated with M.2.

\subsection{Commuting Shed: 0.10}

M.4 receives 0.10 because commuting shed membership is a meaningful but
more indirect community signal than the other primary signals.
Two adjacent tracts may have different primary commuting destinations but
still be a coherent community by land use, housing, and economic character.
The commuting shed signal is most valuable for distinguishing between
adjacent residential tracts that look similar on all physical dimensions
but belong to different functional economic areas.
The 0.10 weight ensures this distinction contributes to the composite
without dominating it.

\subsection{Topographic and Transit: 0.03 and 0.02}

M.5 (topographic community) and M.7 (transit accessibility) receive minimal
weights (3\% and 2\%) for similar reasons:
\begin{itemize}
  \item Their data coverage is more limited than the primary signals
    (GTFS is unavailable for many rural states; SRTM elevation is most
    informative in mountainous states).
  \item Their community type is narrower: topographic community matters
    for redistricting in Colorado, West Virginia, and Alaska but adds
    little in Iowa or Delaware.
  \item Transit accessibility matters for urban redistricting (New York City,
    Chicago, Los Angeles) but is nearly uniform in most of the country.
\end{itemize}
The low weights mean these signals contribute when present without distorting
the composite when absent.
In New York City redistricting, a practitioner might increase M.7 weight
substantially; the \texttt{--community-weights} override supports this.

\subsection{Weight Override Policy}

The default weights represent a general-purpose, jurisdiction-agnostic
calibration.
For specific redistricting proceedings, practitioners should consider:

\begin{itemize}
  \item \textbf{Increase M.6 (zone)} if the state has explicit school district
    preservation criteria or if school district splitting is a hot-button
    issue in the proceeding.
  \item \textbf{Increase M.2 (land use)} if the state has significant water
    boundaries, mountain communities, or sharp urban-rural land use transitions.
  \item \textbf{Increase M.7 (transit)} for city-focused redistricting in
    dense transit cities (New York, Chicago, San Francisco).
  \item \textbf{Decrease M.4 (commuting shed)} for single-employment-center
    states where commuting is uniform (most workers commute to the same CBD).
\end{itemize}

Any weight specification must be declared in advance of the redistricting run
and justified on neutral, non-partisan grounds.
Weight specifications that are chosen to favor specific partisan outcomes
would undermine the legal defensibility of the composite approach.
